\documentclass[12pt,a4paper]{article}
\usepackage[utf8]{inputenc}
\usepackage[spanish]{babel}
\usepackage{geometry}
\usepackage{hyperref}
\usepackage{booktabs}
\usepackage{listings}
\usepackage{xcolor}
\geometry{margin=2.5cm}

% Configuracion de listings para codigo
\lstset{
  basicstyle=\ttfamily\small,
  breaklines=true,
  breakatwhitespace=true,
  frame=single,
  framesep=3pt,
  xleftmargin=5pt,
  xrightmargin=5pt,
  backgroundcolor=\color{gray!10},
  keywordstyle=\color{blue},
  commentstyle=\color{green!50!black},
  stringstyle=\color{red!70!black},
  showstringspaces=false,
  tabsize=2,
  captionpos=b,
  numbers=left,
  numberstyle=\tiny\color{gray},
  numbersep=5pt,
  extendedchars=true,
  literate=
    {á}{{\'{a}}}1
    {é}{{\'{e}}}1
    {í}{{\'{i}}}1
    {ó}{{\'{o}}}1
    {ú}{{\'{u}}}1
    {ñ}{{\~{n}}}1
    {Á}{{\'{A}}}1
    {É}{{\'{E}}}1
    {Í}{{\'{I}}}1
    {Ó}{{\'{O}}}1
    {Ú}{{\'{U}}}1
    {Ñ}{{\~{N}}}1
}

\hypersetup{
  colorlinks=true,
  linkcolor=blue!70!black,
  urlcolor=blue!70!black,
  pdftitle={Evaluacion tecnica del framework forge},
  pdfauthor={Analisis tecnico independiente},
}

\title{\textbf{Evaluación técnica del framework forge}\\
\large para desarrollo con agentes de IA}
\author{Análisis técnico independiente}
\date{3 de mayo de 2026 \\ Versión 1.0}

\begin{document}

\maketitle
\thispagestyle{empty}

\begin{abstract}
El presente informe evalúa \textbf{forge}, un framework de trabajo para equipos que
desarrollan software con agentes de inteligencia artificial. El análisis integra dos
perspectivas técnicas independientes ---una crítica y una favorable--- sobre la misma
base de código, commit \texttt{d828157}, con el objetivo de producir una evaluación
balanceada que permita a equipos de desarrollo tomar una decisión informada sobre la
adopción del framework. La conclusión principal es que forge contiene ideas de diseño
sólidas implementadas parcialmente, cuya aplicabilidad queda limitada a casos de uso
específicos en su estado actual.
\end{abstract}

\tableofcontents
\newpage

% ============================================================
\section{Introducción y contexto}
% ============================================================

\subsection{El problema que forge intenta resolver}

El desarrollo de software asistido por agentes de IA enfrenta un problema de
coordinación que se vuelve crítico a escala de equipo. Sin convenciones explícitas,
los agentes tienden a:

\begin{itemize}
  \item Actuar fuera del scope asignado, modificando componentes no relacionados.
  \item Implementar funcionalidades sin verificación de seguridad (autorización,
        sanitización de inputs, protección de PII).
  \item Producir código inconsistente entre sesiones porque sus reglas están en la
        memoria de cada desarrollador, no en archivos versionados.
  \item Desactualizarse de las prácticas del equipo a medida que el proyecto evoluciona.
\end{itemize}

\textbf{forge} propone una respuesta estructural: codificar las convenciones del equipo
en archivos que los agentes leen, versionados en un repositorio central, y propagados
a los proyectos mediante scripts determinísticos.

\subsection{Descripción del framework}

forge se integra a un proyecto mediante \texttt{git submodule}:

\begin{lstlisting}[language=bash, caption={Instalación básica de forge}]
git submodule add https://github.com/tu-org/forge .agentic
python3 .agentic/scripts/forge-init.py --tool claude-code
\end{lstlisting}

Una vez instalado, el desarrollador describe su proyecto en un \texttt{project.yaml}
y ejecuta \texttt{forge-init.py} para generar la configuración de los agentes.
La estructura del repositorio incluye:

\begin{itemize}
  \item \texttt{core/agents/}: agentes genéricos (orchestrator, backend-engineer,
        security-auditor, compliance-reviewer, test-engineer, frontend-engineer,
        mobile-engineer).
  \item \texttt{profiles/<stack>/agents/}: agentes especializados por stack tecnológico.
  \item \texttt{core/skills/}: skills reutilizables (new-feature, security-audit,
        db-migrate, local2prod, browser-test).
  \item \texttt{adapters/}: scripts para generar configuración en distintos runtimes
        (Claude Code, OpenCode, Kiro).
  \item \texttt{scripts/}: herramientas de gestión (forge-init.py, forge-audit.py).
  \item \texttt{hooks/}: hook de pre-commit para estadísticas de tokens.
\end{itemize}

\subsection{Alcance del análisis}

El análisis cubre el commit \texttt{d828157} del repositorio forge. Los archivos
analizados incluyen: \texttt{README.md}, \texttt{templates/project.yaml.tpl},
\texttt{scripts/forge-init.py}, \texttt{scripts/forge-audit.py},
\texttt{adapters/claude-code/generate-claude-md.py}, siete agentes del core,
\texttt{core/skills/README.md}, cuatro skills, cuatro profiles, \texttt{hooks/pre-commit},
\texttt{docs/agent-standard.md}, y \texttt{core/workflows/sdd.md}.

% ============================================================
\section{Metodología de análisis}
% ============================================================

El análisis se realizó mediante lectura directa del código fuente, sin ejecutar
el framework en un entorno de prueba. Se adoptó un enfoque adversarial estructurado:
dos analistas independientes revisaron el mismo repositorio con instrucciones opuestas
---uno buscando argumentos en contra, otro en favor--- y produjeron informes separados.
El presente documento integra ambos informes con el objetivo de producir una evaluación
libre de sesgos de confirmación.

La evaluación de cada componente siguió tres preguntas:

\begin{enumerate}
  \item ¿Qué problema declara resolver este componente?
  \item ¿El código implementa efectivamente esa solución?
  \item ¿Qué condiciones deben cumplirse para que la implementación aporte valor neto positivo?
\end{enumerate}

Las evidencias citadas son fragmentos directos del código fuente del repositorio,
no inferencias sobre la intención de los autores.

% ============================================================
\section{Hallazgos del análisis crítico}
% ============================================================

\subsection{Los adapters para runtimes alternativos no existen}

El README de forge promete compatibilidad con Claude Code, OpenCode, Codex y otros
runtimes. La inspección del código revela una brecha significativa:

\begin{itemize}
  \item \texttt{adapters/claude-code/}: funcional.
  \item \texttt{adapters/opencode/}: directorio vacío. Cero archivos.
  \item \texttt{adapters/kiro/}: directorio vacío. Cero archivos.
\end{itemize}

La función \texttt{init\_kiro()} en \texttt{forge-init.py} existe y genera dos archivos
\texttt{.md} con contenido genérico. La comparación de complejidad es ilustrativa:
\texttt{init\_claude\_code} tiene aproximadamente 130 líneas de lógica; \texttt{init\_kiro}
tiene aproximadamente 30 y produce documentos de placeholder sin aprovechar ningún
concepto del framework.

\subsection{El ecosistema de profiles cubre cuatro stacks}

El sistema de profiles es la parte conceptualmente más valiosa del framework.
La realidad del inventario en el commit analizado se muestra en la Tabla~\ref{tab:profiles}.

\begin{table}[h]
\centering
\caption{Profiles implementados en forge (mayo 2026)}
\label{tab:profiles}
\begin{tabular}{@{}ll@{}}
\toprule
\textbf{Profile} & \textbf{Agente incluido} \\
\midrule
expo & mobile-engineer.md \\
hono-drizzle & api-engineer.md \\
nextjs-admin & admin-engineer.md \\
playwright-crawler & scanner-engineer.md \\
\midrule
\multicolumn{2}{l}{\textit{Stacks sin profile: Django, Express, Laravel, NestJS, Nuxt}} \\
\bottomrule
\end{tabular}
\end{table}

\subsection{Bug en la función \texttt{install\_agent}}

En \texttt{scripts/forge-init.py}, la función de instalación de agentes contiene
un defecto lógico:

\begin{lstlisting}[language=Python, caption={Bug en install\_agent: status "OK" nunca retorna}]
def install_agent(src: Path, dst: Path, name: str, source_label: str) -> str:
    if not src.exists():
        return "MISS"
    if dst.exists() and not FORCE:
        return "KEEP"
    shutil.copy2(src, dst)
    # Bug: dst.exists() siempre True despues de copy2
    return "UPDATE" if dst.exists() else "OK"
\end{lstlisting}

La evaluación \texttt{dst.exists()} ocurre después de que \texttt{shutil.copy2} ya
copió el archivo. En consecuencia, \texttt{dst.exists()} siempre es \texttt{True}
en ese punto, y el status \texttt{"OK"} nunca se devuelve.

\subsection{El audit recomienda un comando inexistente}

\texttt{forge-audit.py} genera mensajes de corrección de la forma:

\begin{lstlisting}[language=Python, caption={Comando de fix con flag inexistente}]
"fix": f"forge-init.py --tool claude-code --force --only={agent['name']}"
# El flag --only no existe en forge-init.py
\end{lstlisting}

Un desarrollador que siga la recomendación del audit ejecutará un comando que falla.

\subsection{El generador de \texttt{CLAUDE.md} desconecta la fuente de verdad}

El \texttt{project.yaml} contiene una sección \texttt{sprint.phases}. El
\texttt{generate-claude-md.py} genera una sección "Phases activas y estado" con
contenido fijo, no conectado con los datos del YAML. El usuario debe actualizar
manualmente la sección más dinámica del documento más importante que los agentes leerán.

\subsection{El pre-commit hook muta archivos sin visibilidad del desarrollador}

El hook ejecuta \texttt{token-stats.py}, modifica \texttt{docs/progress.html},
y hace \texttt{git add} sobre ese archivo antes de completar el commit. El archivo
entra al commit sin haber sido inspeccionado. Los errores del script se suprimen
con \texttt{|| true}, haciendo que los fallos sean silenciosos.

% ============================================================
\section{Hallazgos del análisis favorable}
% ============================================================

\subsection{La taxonomía de tres tiers resuelve el problema de los agentes God-object}

El \texttt{docs/agent-standard.md} define tres niveles con criterios de clasificación
claros y aplicables:

\begin{itemize}
  \item \textbf{Tier 1 (Universal):} definidos por tipo de output, no por tecnología.
        El orchestrator, security-auditor, compliance-reviewer y test-engineer son
        genuinamente agnósticos al stack.
  \item \textbf{Tier 2 (Profile):} mismo rol que Tier 1 con instrucciones específicas
        al stack. El \texttt{api-engineer} de \texttt{hono-drizzle} sabe que las
        migraciones reversibles requieren un \texttt{down} manual y que el runtime
        en dev es Bun pero en producción es Node 22 LTS.
  \item \textbf{Tier 3 (Dominio):} agentes que conocen conceptos del negocio. Viven
        en el proyecto, no en forge. El framework los reconoce y los verifica en la
        auditoría sin prescribirlos.
\end{itemize}

\subsection{Las reglas de seguridad son específicas y verificables}

Los agentes de forge no mencionan seguridad como aspiración: la codifican como reglas
específicas. El \texttt{backend-engineer} incluye instrucciones no ambiguas:

\begin{lstlisting}[caption={Reglas de seguridad en backend-engineer.md}]
- Usa parametros preparados siempre -- nunca concatenar inputs en queries SQL
- Verifica autenticacion Y autorizacion en cada endpoint
- No loguear PII. Solo IDs hash o indicadores
- No exponer detalles tecnicos de errores al cliente en produccion
\end{lstlisting}

El \texttt{api-engineer} de \texttt{hono-drizzle} agrega restricciones específicas
al stack:

\begin{lstlisting}[caption={Reglas específicas de stack en api-engineer.md}]
- Logs de auditoria son append-only. NUNCA UPDATE ni DELETE
- Multi-tenancy: toda query filtra por tenant_id. Sin excepciones
- Migraciones reversibles: toda migracion tiene down
\end{lstlisting}

\subsection{El skill \texttt{new-feature} codifica un pipeline de calidad completo}

El skill \texttt{core/skills/new-feature/SKILL.md} define seis fases secuenciales
con condiciones de corte: (1) verificar spec aprobada, (2) leer documentación del
área, (3) evaluar si justifica un team de agentes, (4) checklist de seguridad,
(5) orden de implementación (schema $\rightarrow$ tipos $\rightarrow$ backend
$\rightarrow$ frontend $\rightarrow$ build check), y (6) post-implementación completa.

Los agentes no pueden \textit{terminar} una tarea sin haber pasado por security audit
y haber verificado el deploy.

\subsection{El diseño de skills evita el acoplamiento circular}

El \texttt{README.md} de skills documenta explícitamente el grafo de dependencias:

\begin{lstlisting}[caption={Grafo de dependencias de skills}]
new-feature orquesta:
  spec, phase-kickoff, security-audit, db-migrate,
  browser-test (opt), local2prod
local2prod invoca: obsidian-sync (opt)
security-audit: standalone
db-migrate: standalone
\end{lstlisting}

Las dependencias opcionales se resuelven por configuración en \texttt{project.yaml}.
Si \texttt{browser-test} no está en \texttt{skills.active}, \texttt{new-feature}
lo saltea silenciosamente sin romper el pipeline.

\subsection{La auditoría automatizada tiene valor real en CI}

\texttt{forge-audit.py} acepta \texttt{--json} para integración en CI. Detecta:
frontmatter incompleto en agentes, modelo incorrecto (Sonnet en lugar de Opus para
compliance y security), agentes desactualizados respecto al core, agentes huérfanos
no declarados en \texttt{project.yaml}, y oportunidades de activar profiles o skills
disponibles. Esta capacidad transforma los estándares de documentación aspiracional
a reglas ejecutables.

% ============================================================
\section{Análisis comparativo}
% ============================================================

\subsection{Diseño versus implementación}

El análisis revela una brecha consistente entre la calidad del diseño conceptual y
el estado de la implementación. La Tabla~\ref{tab:comparativa} resume los hallazgos
por dimensión.

\begin{table}[h]
\centering
\caption{Análisis comparativo por dimensión}
\label{tab:comparativa}
\begin{tabular}{@{}p{3.5cm}p{5cm}p{5cm}@{}}
\toprule
\textbf{Dimensión} & \textbf{A favor} & \textbf{En contra} \\
\midrule
Arquitectura & Taxonomía 3 tiers coherente & Solo 4 profiles implementados \\
\midrule
Fuente de verdad & project.yaml centraliza todo & Fases del sprint no se propagan \\
\midrule
Seguridad & Reglas específicas y verificables & Compliance sin acceso a texto legal \\
\midrule
Multi-runtime & Diseño limpio de separación & OpenCode y Kiro: vacíos \\
\midrule
Auditoría & Output JSON integrable en CI & Flag --only del fix no existe \\
\midrule
Actualización & Propagación desde repo central & --force destruye customizaciones \\
\midrule
Adopción & Setup en un solo comando & Sin comando de teardown \\
\bottomrule
\end{tabular}
\end{table}

\subsection{Comparación con alternativas}

La alternativa más común es escribir agentes directamente en \texttt{.claude/agents/}
sin framework: costo de 30 minutos para un \texttt{CLAUDE.md} inicial, control total,
sin dependencias. forge agrega valor cuando el equipo necesita: (a) propagación
consistente de reglas de seguridad entre proyectos, (b) agentes especializados por
stack mantenidos centralmente, (c) auditoría en CI, o (d) compliance con poder de veto.

\subsection{Riesgos sistémicos}

Los riesgos más significativos no son los bugs individuales sino los patrones sistémicos:

\begin{itemize}
  \item \textbf{Actualización destructiva:} el mecanismo \texttt{--force} sobreescribe
        customizaciones sin merge. Con el tiempo, los equipos dejan de actualizar los
        agentes para proteger sus customizaciones.
  \item \textbf{Dependencia sin escape:} no existe comando de teardown. La salida
        requiere limpieza manual que puede ser omitida.
  \item \textbf{Falsa seguridad en compliance:} el \texttt{compliance-reviewer} es
        un primer filtro valioso, pero el framework no advierte que opera sin acceso
        a texto legal oficial.
\end{itemize}

% ============================================================
\section{Conclusiones y recomendaciones}
% ============================================================

\subsection{Conclusión general}

forge es un framework con arquitectura conceptual correcta e implementación en estado
temprano. La taxonomía de tres tiers resuelve un problema real de forma elegante.
Las reglas de seguridad codificadas en los agentes son específicas, verificables y
no se degradan entre sesiones. El pipeline del skill \texttt{new-feature} define una
barra de calidad ejecutable. Sin embargo, los adapters para runtimes alternativos son
directorios vacíos, el ecosistema de profiles cubre cuatro stacks, y el tooling de
auditoría tiene bugs que afectan directamente la experiencia del usuario.

\subsection{Condiciones de adopción recomendable}

Se recomienda adoptar forge cuando se cumplen simultáneamente:

\begin{enumerate}
  \item Equipo de 3 a 8 personas usando agentes de IA activamente en el loop de desarrollo.
  \item Stack cubierto por alguno de los cuatro profiles existentes.
  \item Claude Code como runtime principal.
  \item Necesidades de compliance reales (GDPR, Ley 21.719, LGPD).
  \item Tolerancia para customizaciones mínimas de agentes.
\end{enumerate}

\subsection{Recomendaciones para equipos adoptantes}

\begin{itemize}
  \item Documentar la política de uso de \texttt{--force}: cuándo está permitido
        sobreescribir agentes customizados.
  \item No confiar en los mensajes de fix del audit para comandos con el flag
        \texttt{--only}: ese flag no existe en \texttt{forge-init.py}.
  \item Verificar manualmente la sección de fases del sprint en \texttt{CLAUDE.md}
        después de cada regeneración.
  \item Tratar el \texttt{compliance-reviewer} como primer filtro, no como revisión
        legal suficiente.
\end{itemize}

\subsection{Recomendaciones para el equipo de forge}

\begin{itemize}
  \item Implementar los adapters para OpenCode y Kiro o remover su mención del
        README hasta que estén disponibles.
  \item Corregir el flag \texttt{--only} en los mensajes del audit o implementar
        el flag en \texttt{forge-init.py}.
  \item Conectar la sección \texttt{sprint.phases} de \texttt{project.yaml} con
        la generación de \texttt{CLAUDE.md}.
  \item Agregar un comando de teardown documentado.
  \item Implementar un mecanismo de merge para actualizaciones de agentes que no
        destruya customizaciones.
\end{itemize}

% ============================================================
\appendix
\section{Estado de los adapters (mayo 2026)}
% ============================================================

\begin{lstlisting}[caption={Estado del directorio adapters/}]
adapters/
|-- claude-code/
|   |-- generate-claude-md.py   <- implementado
|   +-- commands/wiki/          <- implementado
|-- opencode/                   <- directorio vacio
+-- kiro/                       <- directorio vacio
\end{lstlisting}

% ============================================================
\section{Bug completo en \texttt{install\_agent}}
% ============================================================

\begin{lstlisting}[language=Python, caption={Función install\_agent completa con el bug}]
def install_agent(src: Path, dst: Path, name: str, source_label: str) -> str:
    if not src.exists():
        return "MISS"
    if dst.exists() and not FORCE:
        return "KEEP"
    shutil.copy2(src, dst)
    # Bug: dst.exists() siempre es True despues de copy2
    # El status "OK" para nueva instalacion nunca se retorna
    return "UPDATE" if dst.exists() else "OK"
\end{lstlisting}

% ============================================================
\section{Grafo de dependencias de skills}
% ============================================================

\begin{lstlisting}[caption={Grafo de dependencias documentado en core/skills/README.md}]
new-feature orquesta:
  spec, phase-kickoff, security-audit, db-migrate,
  browser-test (opt), local2prod

local2prod invoca:
  obsidian-sync (opt)

security-audit: standalone
db-migrate: standalone
\end{lstlisting}

\bigskip
\hrule
\bigskip
\noindent\textit{Informe generado el 3 de mayo de 2026.
Base de análisis: commit \texttt{d828157} del repositorio forge.}

\end{document}
